\chapter{正文元素示例}

\section{公式示例}

设接收信号为 $y=x+n$，其中 $n$ 为加性高斯白噪声，则对数似然比可写为

\begin{equation}
LLR = \frac{2y}{\sigma^2}
\end{equation}

其中，$y$ 表示接收符号，$\sigma^2$ 表示噪声方差。该类公式可用于演示编号、对齐以及文中引用方式。

\section{表格示例}

表\ref{tab:sample_params} 给出了一个简单的参数表示例。

\begin{bupttable}{示例参数表}{tab:sample_params}
    \begin{tabular}{@{}p{3cm} p{3cm} p{6cm}@{}}
        \toprule
        参数名称 & 取值 & 说明 \\
        \midrule
        码长 $N$ & 1024 & 示例中使用的块长度 \\
        码率 $R$ & 0.5 & 示例中使用的码率 \\
        信道模型 & AWGN & 加性高斯白噪声信道 \\
        \bottomrule
    \end{tabular}
\end{bupttable}

\section{图示例}

图\ref{fig:sample_flow} 展示了一个简单的流程示意图。

\begin{figure}[!htbp]
    \centering
    \begin{tikzpicture}[
        node distance=5mm and 8mm,
        box/.style={draw, rounded corners=2pt, align=center, minimum width=2.4cm, minimum height=0.9cm, font=\wuhao},
        line/.style={-{Latex[length=2mm]}, thick}
    ]
        \node[box] (a) {输入};
        \node[box, right=of a] (b) {处理};
        \node[box, right=of b] (c) {输出};
        \draw[line] (a) -- (b);
        \draw[line] (b) -- (c);
    \end{tikzpicture}
    \caption{示例流程图}
    \label{fig:sample_flow}
\end{figure}

\section{本章小结}

本章展示了公式、表格和图在模板中的基本书写方式，可作为后续撰写实际论文时的参考。
